%!TEX encoding = UTF-8 Unicode
%!TEX root = reporte-derecho-penal-especial-mexicano-Actividad-8.tex
% arara: xelatex: { interaction: nonstopmode, synctex: true }
% arara: biber
% arara: xelatex: { interaction: nonstopmode, synctex: true }
% arara: xelatex: { interaction: nonstopmode, synctex: true }

\documentclass[12pt, a4paper, oneside]{article}
\usepackage[aulatex]{template-unadm}
\usepackage{template-unadm-ld}
\usepackage{template-unadm-lex}
\usepackage{pdflscape}
\usepackage{longtable}
\usepackage{array}
\newcolumntype{L}[1]{>{\raggedright\let\newline\\\arraybackslash\hspace{0pt}}m{#1}}
\newcolumntype{C}[1]{>{\centering\let\newline\\\arraybackslash\hspace{0pt}}m{#1}}
\newcolumntype{R}[1]{>{\raggedleft\let\newline\\\arraybackslash\hspace{0pt}}m{#1}}

% Datos canónicos del estudiante
\def\documentauthor{Mart\'in Jonathan de la Cruz Mu\~noz}
\author{Mart\'in Jonathan de la Cruz Mu\~noz}
\def\documentauthorcanonico{Martín Jonathan de la Cruz Muñoz}  % UTF-8 para validación
\def\documentsubject{Licenciatura en Derecho}
\def\documenttitle{Juego de roles: Delitos contra la paz y seguridad de las personas}
\def\documenttitlehf{Juego de roles: Delitos contra la paz}
\def\documentsubtitle{Análisis jurídico de amenazas, cobranza ilegítima, allanamiento de morada y usurpación de identidad}
\def\documentdate{Semana 8, 2026}
\def\documentperiodo{2026-2 B2}
\def\documentcareer{Licenciatura en Derecho}
\def\documentinstitution{UnADM}
\def\documentmatricula{ES2611202040}
\def\documentgroup{020}
\def\documentcourse{Derecho Penal Especial Mexicano}
\def\documentteacher{}

\hfwidthwrap{false}
\hfwidthtitle{0.55}
\hfwidthcourse{0.42}

\begin{document}
% Wrapper local para apuntar a la plantilla maestra de AulaTeX
\input{C:/Users/Sysx/Documents/AulaTeX-Academico/aulaTEX/base/Plantilla-Informe/template.tex}


\begin{abstractd}
El presente documento desarrolla la Actividad S8 de Derecho Penal Especial Mexicano, consistente en un juego de roles sobre delitos contra la paz, la seguridad de las personas y la inviolabilidad del domicilio. Se analizan cuatro figuras delictivas: amenazas, cobranza ilegítima, allanamiento de morada y usurpación de identidad, presentando una tabla comparativa de sus elementos normativos y tres consultas jurídicas simuladas como asesor penal.
\end{abstractd}

\templateIndex

\section{Marco normativo de los delitos contra la paz y seguridad}

Los delitos contra la paz, la seguridad de las personas y la inviolabilidad del domicilio se encuentran regulados en el Código Penal Federal (CPF) y en los códigos penales de las entidades federativas. Estos delitos protegen bienes jurídicos fundamentales como la tranquilidad personal, la seguridad individual y el derecho a la intimidad.

El análisis de estos delitos requiere una perspectiva garantista que equilibre la protección de los bienes jurídicos con el respeto a los derechos fundamentales de los imputados. La taxatividad penal y el principio de ultima ratio deben guiar la interpretación de estos tipos penales para evitar arbitrariedades en su aplicación.

\section{Tabla comparativa de figuras delictivas}

\clearpage
\pagestyle{empty}
\begin{landscape}
\begin{center}
\scriptsize
\renewcommand{\arraystretch}{1.2}
\setlength{\tabcolsep}{4pt}
\begin{longtable}{|L{3cm}|L{3cm}|L{4cm}|L{4cm}|L{4cm}|}
\hline
\textbf{Figura delictiva} & \textbf{Ley/Artículo} & \textbf{Conducta prevista} & \textbf{Elementos a comprobar} & \textbf{Observaciones} \\ \hline

Amenazas & CPF Art. 282 & Amenazar a otro con causarle un mal en su persona, bienes, honor o derechos, o en la persona, bienes, honor o derechos de alguien con quien esté ligado por algún vínculo. & 1. Existencia de una amenaza; 2. Que el mal anunciado sea futuro, posible y dependiente de la voluntad del amenazante; 3. Que la amenaza sea seria y verosímil; 4. Que cause alarma o miedo en la víctima. & La amenaza debe ser idónea para causar alarma o miedo en una persona razonable. No se requiere que el amenazante tenga la intención real de ejecutar el mal anunciado. \\ \hline

Cobranza ilegítima & CPF Art. 282 Bis & Realizar actos de cobranza con violencia, intimidación, amenazas o cualquier otro medio ilícito para exigir el pago de una deuda. & 1. Existencia de una deuda; 2. Actos de cobranza realizados por el acreedor o un tercero; 3. Uso de violencia, intimidación, amenazas o cualquier otro medio ilícito; 4. Intención de exigir el pago de la deuda. & Este delito protege la tranquilidad y seguridad de las personas frente a prácticas abusivas de cobranza. La violencia o intimidación puede ser física o psicológica. \\ \hline

Allanamiento de morada & CPF Art. 285 & Entrar en morada ajena sin permiso del titular o contra su voluntad, o mantenerse en ella después de haber sido requerido para que se retire. & 1. Entrada en morada ajena; 2. Falta de consentimiento del titular; 3. Intención de entrar o mantenerse en la morada sin derecho. & La morada comprende no solo la vivienda, sino también sus dependencias cerradas. El consentimiento del titular excluye la tipicidad del delito. \\ \hline

Usurpación de identidad & CPF Art. 211 Bis & Apoderarse, transferir, utilizar o divulgar, sin consentimiento, datos personales de otra persona, con el fin de obtener un lucro indebido o causar un perjuicio. & 1. Apoderamiento, transferencia, utilización o divulgación de datos personales; 2. Falta de consentimiento del titular; 3. Fin de obtener un lucro indebido o causar un perjuicio. & Este delito protege el derecho a la identidad y la intimidad de las personas. Los datos personales incluyen cualquier información concerniente a una persona física identificada o identificable. \\ \hline

\end{longtable}
\end{center}
\end{landscape}
\clearpage
\pagestyle{fancy}

\section{Consultas jurídicas simuladas}

\subsection{Caso 1: Amenazas en contexto de violencia familiar}

**Consulta:** Una mujer recibe mensajes de texto de su expareja donde le advierte que "si no regresa con él, le hará daño a sus hijos". La mujer acude a usted como asesor jurídico. ¿Qué acciones legales recomendaría?

**Respuesta:** Recomendaría a la mujer presentar una denuncia penal por el delito de amenazas (Art. 282 CPF) y solicitar medidas de protección conforme a la Ley General de Acceso de las Mujeres a una Vida Libre de Violencia. Además, sugeriría iniciar un procedimiento de violencia familiar y solicitar la intervención de las autoridades competentes para garantizar la seguridad de ella y sus hijos.

El delito de amenazas se configura porque:
1. Existe una amenaza clara y concreta;
2. El mal anunciado (daño a los hijos) es futuro, posible y dependiente de la voluntad del amenazante;
3. La amenaza es seria y verosímil, considerando el contexto de violencia familiar;
4. La amenaza causa alarma y miedo en la víctima.

\subsection{Caso 2: Cobranza ilegítima con prácticas abusivas}

**Consulta:** Un despacho de cobranza visita el domicilio de un deudor y le exige el pago de una deuda vencida. Durante la visita, los cobradores gritan, insultan al deudor y amenazan con "hacerle la vida imposible" si no paga. ¿Qué delito se configura y qué acciones legales proceden?

**Respuesta:** Se configura el delito de cobranza ilegítima (Art. 282 Bis CPF) debido a que los cobradores utilizaron intimidación y amenazas para exigir el pago de la deuda. Recomendaría al deudor presentar una denuncia penal y solicitar la intervención de la Procuraduría Federal del Consumidor (PROFECO) para que investigue las prácticas abusivas del despacho de cobranza.

Los elementos del delito son:
1. Existencia de una deuda;
2. Actos de cobranza realizados por un tercero;
3. Uso de intimidación y amenazas;
4. Intención de exigir el pago de la deuda mediante prácticas ilícitas.

\subsection{Caso 3: Allanamiento de morada en contexto de conflicto vecinal}

**Consulta:** Un vecino ingresa al patio trasero de una casa sin permiso del propietario para retirar unas macetas que, según él, están invadiendo su propiedad. El propietario lo descubre y lo denuncia por allanamiento de morada. ¿Procede la denuncia?

**Respuesta:** La denuncia por allanamiento de morada (Art. 285 CPF) procede porque el vecino ingresó a una dependencia cerrada de la morada ajena sin consentimiento del titular. Aunque el motivo del ingreso fue retirar unas macetas, este hecho no justifica la violación a la inviolabilidad del domicilio.

Los elementos del delito son:
1. Entrada en morada ajena (patio trasero);
2. Falta de consentimiento del titular;
3. Intención de ingresar sin derecho.

Recomendaría al propietario presentar la denuncia penal y, en paralelo, iniciar un procedimiento civil para resolver el conflicto de límites entre las propiedades.

\section{Conclusiones}
\clearpage

Considero que los delitos contra la paz, la seguridad de las personas y la inviolabilidad del domicilio protegen bienes jurídicos fundamentales que son esenciales para la convivencia social y el desarrollo individual. La taxatividad penal y el principio de ultima ratio deben guiar la interpretación de estos tipos penales para evitar arbitrariedades en su aplicación y garantizar que solo las conductas más graves sean sancionadas penalmente.

En el caso de las amenazas, sostengo con firmeza que este delito debe interpretarse bajo una perspectiva de género, especialmente en contextos de violencia familiar. Las amenazas en estos casos no solo afectan la tranquilidad de la víctima, sino que también generan un clima de terror que limita su libertad y autonomía.

Respecto a la cobranza ilegítima, rechazo frontalmente las prácticas abusivas utilizadas por algunos despachos de cobranza. Estas conductas no solo vulneran la tranquilidad y seguridad de las personas, sino que también generan un clima de impunidad que afecta la confianza en las instituciones. Es fundamental que las autoridades competentes, como la PROFECO, actúen con mayor firmeza para sancionar estas prácticas y proteger los derechos de los deudores.

En cuanto al allanamiento de morada, es crucial distinguir entre los conflictos civiles y las conductas delictivas. La inviolabilidad del domicilio es un derecho fundamental que debe ser protegido con firmeza, pero también es importante que los conflictos vecinales se resuelvan a través de mecanismos alternativos de solución de controversias para evitar la judicialización innecesaria de estos casos.

\footnote{Declaración: Este trabajo académico fue desarrollado por \documentauthorcanonico\ (matrícula \documentmatricula) como parte de las actividades de la asignatura \documentcourse\ de la \documentsubject\ en la \documentinstitution. Se utilizó apoyo tecnológico para la organización y redacción, pero el contenido, análisis y conclusiones son responsabilidad exclusiva del autor, conforme a los lineamientos institucionales de la UnADM sobre el uso de herramientas de asistencia académica.}

\clearpage
\printbibliography
\end{document}